%========================================================================================
%   Chapter 5 -- Implementation (merged from the former chapters 5 and 6)
%   Code listings were removed to shorten the report.
%========================================================================================
\chapter{Implementation}\label{chap:implementation}
\label{ch:implementation}

This chapter describes how the architecture and the pipeline designed in Chapter~\ref{chap:design} were realized, without restating the design: the development environment (Section~\ref{sec:impl-env}), the backend (Section~\ref{sec:impl-back}), the frontend (Section~\ref{sec:impl-front}), the container images (Section~\ref{sec:impl-docker}), Terraform (Section~\ref{sec:impl-terraform}), Ansible (Section~\ref{sec:impl-ansible}) and the GitHub Actions workflows (Section~\ref{sec:impl-cicd}). Measurements and scan results are presented in Chapter~\ref{chap:results}.

%========================================================================================
\section{Development Environment and Repository Layout}
\label{sec:impl-env}

The backend is a Java~21 application built with Maven (wrapper \texttt{mvnw}) on Spring Boot~3.5.16 (3.5.6 at the time of the measurements of Chapter~\ref{chap:results}). The main dependencies are Spring Web, Spring Data JPA with the PostgreSQL driver, Spring Security, Bean Validation, \texttt{jjwt}~0.12.5, Lombok, and \texttt{springdoc-openapi-starter-webmvc-ui}~2.8.9; the Actuator starter exposes the health endpoint, H2 and \texttt{spring-security-test} are used for tests, and the JaCoCo plugin (0.8.12) produces coverage reports. The frontend is a Create React App project using React~19.1, Material~UI~7.3, React Router~7 and Axios. The database is PostgreSQL~16. The whole system lives in a single repository organized by concern (\texttt{back/}, \texttt{front/}, \texttt{infra/}, \texttt{ansible/}, \texttt{docker/}, \texttt{.github/workflows/}); the directory \texttt{docs/} holds this report.

Runtime configuration is injected through environment variables, never committed. The file \texttt{.env.example} documents the seven variables the application needs (\texttt{POSTGRES\_DB}, \texttt{POSTGRES\_USER}, \texttt{POSTGRES\_PASSWORD}, \texttt{JWT\_SECRET}, \texttt{APP\_USER\_USERNAME}, \texttt{APP\_USER\_PASSWORD}, \texttt{APP\_CORS\_ALLOWED\_ORIGINS}); developers copy it to a git-ignored \texttt{.env}. Spring configuration is split into a common file, which sets the application name and exposes only the health endpoint, and two profiles: \texttt{local} points the datasource at \texttt{localhost:5432} and lets Hibernate update the schema, while \texttt{docker} defines the application user and token settings from the environment, the datasource being supplied by Compose. The token lifetime is 1\,800\,000~ms (30 minutes) in both. No further profiles exist.

Two ways of running the system locally are provided: \texttt{back/start-local.sh} loads \texttt{.env} and runs \texttt{./mvnw spring-boot:run} with the \texttt{local} profile, and \texttt{start.sh} at the repository root rebuilds and starts the three containers of \texttt{docker/docker-compose.local.yaml} (Section~\ref{sec:impl-docker}).

%========================================================================================
\section{Backend: Spring Boot REST API}
\label{sec:impl-back}

\subsection{Domain, Persistence and Services}

The package structure follows Figure~\ref{fig:backend-layers}; the constants of \texttt{ApiPaths} compose the versioned prefix \texttt{/api/v1}. Three JPA entities map the domain. A \texttt{TransportLog} references one client and one driver and stores the load and unload dates and locations, the destination, the delivery note, the monetary and fuel figures of the trip (advance, fuel quantity and price per liter, variable charge, charge price, client tariff and trip price), and the operator and commercial responsible for it. Lombok generates accessors and builders. The three repositories simply extend \texttt{JpaRepository}, so all queries are derived and parameterized and no SQL string is concatenated anywhere in the code.

The API never exposes entities: every module has a creation request, an update request and a response, declared as immutable Java records with Jakarta Bean Validation constraints, plus the generic \texttt{BaseResponse<T>} and \texttt{PagedResponse<T>}. One \texttt{@Component} mapper per module converts between records and entities; explicit mappers were preferred to MapStruct since, with three small entities, they are easy to read and debug. The transport-log mapper receives the already loaded client and driver and delegates the nested conversions to the two other mappers.

Each service is an interface with an implementation receiving its collaborators by constructor injection. The implementations load by identifier and raise a \texttt{ResourceNotFoundException} when the entity does not exist (also on delete), update, and build the paginated answer from a Spring Data \texttt{Page}. The transport-log service is the only one spanning several aggregates: its create and update operations are \texttt{@Transactional} and first resolve the referenced client and driver, so that a request naming an unknown one fails with a not-found error before anything is written.

\subsection{Controllers and OpenAPI Documentation}

The controllers expose the operations of Table~\ref{tab:api}. The \texttt{size} parameter is constrained to 1--100 by Bean Validation on a \texttt{@Validated} controller, and the transport-log path is \texttt{/trans\_logs}. Every controller is annotated for springdoc: \texttt{@Tag} groups the operations, \texttt{@Operation} describes them, and \texttt{@ApiResponse} points to \texttt{BaseResponse}, \texttt{PagedResponse} and \texttt{ErrorResponse}. A class-level \texttt{@SecurityRequirement} (scheme \texttt{BearerAuth}) marks the endpoints as requiring the bearer scheme declared once in \texttt{OpenApiConfig}, where the title, version, license and server are also defined; the login operation overrides it with an empty requirement. The specification is served at \texttt{/v3/api-docs} (JSON) and \texttt{/v3/api-docs.yaml}, and the workflow \texttt{openapi-docs.yml} (Section~\ref{sec:impl-cicd}) publishes a static Swagger~UI built from it to GitHub Pages after each successful backend pipeline on the main branch.

\subsection{Security and Error Handling}
\label{sec:impl-security}

The security design of Chapter~\ref{chap:design} is realized with Spring Security in a stateless configuration, as follows.

\paragraph{Authentication.} \texttt{SimpleUserDetailsService} hashes the password bound from \texttt{app.user.*} with BCrypt once, in a \texttt{@PostConstruct} method, discards the clear text, and returns a \texttt{UserDetails} only for the configured user name. \texttt{@EnableMethodSecurity} is switched on but no method-level rule is declared. The \texttt{AuthController} passes the credentials of a validated \texttt{LoginRequestDto} to the \texttt{AuthenticationManager}; on success the \texttt{JwtTokenProvider} builds a token whose subject is the user name and whose \texttt{roles} claim lists the authorities, signed with an HMAC key derived from \texttt{JWT\_SECRET}, and the response contains the token, the user name and the expiry time.

\paragraph{Filter chain.} The \texttt{JwtAuthenticationFilter}, a \texttt{OncePerRequestFilter} registered before the standard username/password filter, skips the public paths, extracts the bearer token, validates its signature and expiry, loads the user and stores an authentication object in the security context. An invalid or missing token leaves the context empty and the decision is taken by the authorization rules (\texttt{anyRequest().authenticated()} after the few public paths); the filter itself contains no exception for any business path, so all client endpoints are protected. The chain registers an \texttt{authenticationEntryPoint} and an \texttt{accessDeniedHandler} that write the same JSON \texttt{ErrorResponse} with status \texttt{401} or \texttt{403}; without them, Spring Security would answer an unauthenticated request with an empty \texttt{403}, which the frontend could not distinguish from a permission problem. Sessions are disabled and CSRF protection is switched off.

\paragraph{CORS.} \texttt{CorsConfig} registers a \texttt{CorsConfigurationSource} that allows the origins of \texttt{app.cors.allowed-origins} (filled from \texttt{APP\_CORS\_ALLOWED\_ORIGINS}), the usual HTTP methods, any header and credentials, and exposes the \texttt{Authorization} header. The Ansible environment template sets the same variable on the servers.

\paragraph{Error handling.} A \texttt{@ControllerAdvice} class, \texttt{GlobalExceptionHandler}, converts exceptions into the \texttt{ErrorResponse} record: failed Bean Validation yields a 400 listing the offending fields, \texttt{ResourceNotFoundException} a 404, \texttt{BadCredentialsException} a 401 with a deliberately generic message, and any other exception a 500. Failures detected before a request reaches a controller are reported by the security configuration with the same body. The \texttt{@ApiResponse} annotations reference the same type, so the documented and the actual error formats coincide.

\subsection{Test Strategy}

The backend tests follow the layers of the application. The Driver module served as the reference for each level and was reproduced for the Client and TransportLog modules. The suite contains 82 tests, all passing in the last run, which measured a JaCoCo line coverage of 77.1\,\% (see Chapter~\ref{chap:results}).
\begin{itemize}
  \item \textbf{Service unit tests} use JUnit~5 and Mockito, with the repository and the mapper mocked.
  \item \textbf{Web-layer tests} use \texttt{@WebMvcTest} with \texttt{MockMvc} and a \texttt{@MockBean} service; an imported \texttt{TestSecurityConfig} registers the real JWT components but a permissive filter chain, so the tests concentrate on routing, validation and serialization (201, 200, 400, 404, and the 401 answer to bad credentials for the \texttt{AuthController}).
  \item \textbf{Repository tests} use \texttt{@DataJpaTest} against an in-memory H2 database.
  \item \textbf{Integration tests} use \texttt{@SpringBootTest} with \texttt{@AutoConfigureMockMvc}; environment properties are provided with \texttt{@TestPropertySource}, an authenticated user is simulated with \texttt{@WithMockUser}, and the table is emptied before each test.
  \item \textbf{Security integration tests} (\texttt{SecurityIntegrationTest}) load the real filter chain and check that a missing, malformed or expired token on a business endpoint (including the client endpoints) and a wrong password yield \texttt{401}, and that a valid login followed by a request with the issued token gives access.
\end{itemize}
JaCoCo instruments the test run, and the pipeline uploads its report to Codecov (Section~\ref{sec:impl-cicd}).

%========================================================================================
\section{Frontend: React Single-Page Application}
\label{sec:impl-front}

The source tree of \texttt{front/src} has a \texttt{components} folder with one sub-folder per screen (\texttt{Auth}, \texttt{Dashboard}, \texttt{Clients}, \texttt{Drivers}, \texttt{TransportLogs}), a \texttt{Layout} component (application bar, navigation menu, logout) and a \texttt{Common} folder with \texttt{DataTable}, \texttt{ConfirmDialog}, \texttt{LoadingSpinner} and \texttt{ErrorBoundary}; \texttt{contexts} holds the \texttt{AuthContext}, \texttt{services} the Axios client, and \texttt{utils} some formatting helpers.

In routing, the public \texttt{/login} route is wrapped in a \texttt{PublicRoute} that redirects an already authenticated user to the dashboard; all other screens are children of a route rendering the \texttt{Layout} inside a \texttt{ProtectedRoute}, which shows a spinner while the stored session is restored and redirects anonymous visitors to the login page. Unknown paths fall back to the dashboard. This protection is a usability feature; the authoritative checks are performed by the backend.

The \texttt{AuthContext} exposes \texttt{isAuthenticated}, \texttt{user}, \texttt{loading}, \texttt{login} and \texttt{logout} through a provider and a \texttt{useAuth} hook. The login function calls the backend, validates the shape of the answer and only then marks the session as authenticated; there is no local or simulated login path. The token and user name are written to \texttt{localStorage} and restored at application start; this survives page reloads but makes the token readable by any script in the page, a trade-off discussed with the limitations.

All backend calls go through one Axios instance whose methods mirror the endpoints of Table~\ref{tab:api}. Its response interceptor stores user-facing messages in \texttt{error.userMessage} (time-out, network error, 403, 5xx, or the \texttt{message} returned by the API). The base URL is \texttt{process.env.REACT\_APP\_API\_URL} when set at build time and the relative path \texttt{/api/v1} otherwise: in the containers nginx proxies \texttt{/api/} to the backend, and in local development the \texttt{proxy} entry of \texttt{package.json} forwards unknown requests to \texttt{http://localhost:8080}. The same build can therefore be deployed to any host without recompilation.

The \textbf{Dashboard} issues three requests in parallel with \texttt{Promise.all}: it reads \texttt{totalElements} of a one-element page of clients and of drivers for the counts, and fetches the first five transport logs. The \textbf{Clients} and \textbf{Drivers} screens share one pattern: a paginated \texttt{DataTable}, a dialog with the creation and edition form and field-level validation, a \texttt{ConfirmDialog} before deletion, and a \texttt{Snackbar} for feedback. The \textbf{Transport logs} screen uses \texttt{Autocomplete} fields fed by the client and driver endpoints, and checks mandatory and numeric fields before submission; monetary values, including the trip price, are entered by the operator and not derived by the application. \texttt{DataTable} is driven by a column description and by the pagination state held by the parent, which maps directly onto the server-side pagination; \texttt{ErrorBoundary} shows a recovery screen instead of a blank page on rendering errors. The frontend tests are limited to smoke tests in \texttt{App.test.js}, which the pipeline executes.

In the container image the compiled application is served by nginx. The \texttt{location /} block implements the single-page-application fallback with \texttt{try\_files}, so that deep links such as \texttt{/drivers} are answered by \texttt{index.html}; static assets are cached for a year and HTML never. The \texttt{location /api/} block proxies to the container \texttt{spring-boot-app} on port 8080 with forwarding headers and 30-second time-outs, and Swagger~UI and the OpenAPI document are proxied likewise. The server also enables gzip, adds the headers \texttt{X-Frame-Options}, \texttt{X-Content-Type-Options}, \texttt{X-XSS-Protection}, \texttt{Referrer-Policy} and a \texttt{Content-Security-Policy}, denies hidden files, and answers \texttt{/health} directly for the image health check.

%========================================================================================
\section{Containerization}
\label{sec:impl-docker}

The backend image is built in two stages. The first uses the Temurin JDK~21 image, copies \texttt{pom.xml} and the Maven wrapper before the sources so that dependencies are resolved in a separate cached layer, and packages the application with tests skipped, since the pipeline has already run them. The second starts from the Temurin JRE~21 image, creates a dedicated system user and group, copies only the executable JAR owned by that user and runs as this unprivileged user, so a compromised process does not hold root privileges in the container.

The frontend image also has two stages: Node~20 on Alpine compiles the React application with \texttt{npm ci} and \texttt{npm run build}, and the static files are copied into an \texttt{nginx:alpine} image together with the nginx configuration above. The runtime stage upgrades the Alpine packages, installs \texttt{curl} for the health check, gives the \texttt{nginx} user ownership of the cache, log and content directories and switches to \texttt{USER nginx}; a \texttt{HEALTHCHECK} polls \texttt{/health} every 30 seconds. The file is named \texttt{dockerfile} in lower case.

Three Compose files live in \texttt{docker/}. \texttt{docker-compose.local.yaml} builds the backend and frontend from the sources and starts a \texttt{postgres:16-alpine} database (used by \texttt{start.sh}); the other two use an already built backend image, \texttt{docker-compose.back.yaml} starting only the backend and the database. The services share a bridge network and the database keeps its data in the named volume \texttt{db\_data} (Figure~\ref{fig:docker-arch}). The startup order is guarded by a health check rather than a mere \texttt{depends\_on}: the database declares a \texttt{pg\_isready} probe and the backend waits for \texttt{service\_healthy}, so Spring does not open its connection pool before PostgreSQL accepts connections. Credentials come from \texttt{.env} through variable substitution, and the \texttt{docker} profile is activated through \texttt{SPRING\_PROFILES\_ACTIVE}.

\begin{figure}[htbp]
\centering
\begin{adjustbox}{max height=4.2cm}
\begin{tikzpicture}[node distance=1.6cm, >=latex, thick]

\tikzstyle{container}=[
  rectangle, draw=Eblue, fill=Ecyan!15,
  rounded corners=6pt, minimum width=4.4cm, minimum height=1.2cm,
  font=\small, align=center
]
\tikzstyle{volume}=[
  cylinder, shape border rotate=90,
  draw=Eblue, fill=Eorange!30,
  minimum height=1.2cm, minimum width=1.6cm, aspect=0.3,
  font=\small
]

\node[container](frontend){\textbf{React frontend}\\[-1pt]\scriptsize published on port 3000};
\node[container, below=of frontend](backend){\textbf{Spring Boot backend}\\[-1pt]\scriptsize \texttt{app}, port 8080};
\node[container, below=of backend](database){\textbf{PostgreSQL 16}\\[-1pt]\scriptsize \texttt{db}, port 5432};
\node[volume, right=4.0cm of database](dbvol){\texttt{db\_data}};

\draw[<->, Eblue] (frontend) -- node[right]{\scriptsize HTTP / REST} (backend);
\draw[<->, Eblue] (backend) -- node[right]{\scriptsize JDBC} (database);
\draw[->, Eblue] (database) -- node[above, pos=0.62]{\scriptsize named volume} (dbvol);

\node[draw, dashed, rounded corners=10pt,
      fit=(frontend)(backend)(database),
      inner sep=0.5cm, label={[Eblue,font=\small]below:Compose network \texttt{my-app-network} (bridge)}] {};

\end{tikzpicture}
\end{adjustbox}
\caption{Services of the local Docker Compose stack (\texttt{docker/docker-compose.local.yaml}).}
\label{fig:docker-arch}
\end{figure}


%========================================================================================
\section{Infrastructure as Code with Terraform}
\label{sec:impl-terraform}

The \texttt{infra/} directory is a Terraform root module for the DigitalOcean provider (constraint \texttt{\textasciitilde{}> 2.43}, Terraform 1.5 or later). The state is kept in the Terraform Cloud backend, and the provider token is a sensitive variable supplied by the pipeline. \texttt{main.tf} creates the shared VPC and instantiates \texttt{modules/stack} twice (topology in Figure~\ref{fig:terraform-vpc-architecture}); the environments differ only in the module arguments, mainly droplet sizes with different defaults (a larger frontend and backend for production). Variable validation encodes a security rule in the interface: the list of CIDR blocks allowed to reach the bastion over SSH is rejected if it contains \texttt{0.0.0.0/0}.

The \texttt{stack} module is parameterized by the environment name, region, VPC identifier, SSH keys, application port and machine sizes. It creates three droplets from an Ubuntu~22.04 image, attached to the VPC, monitored and tagged by environment and role; three cloud firewalls bound to them, whose tag-based sources keep the rules valid when addresses change; and a managed PostgreSQL cluster (engine and version from variables, a single node of the smallest size, private network, Monday maintenance window) with a database named after the environment and a database firewall whose only rule admits the backend droplet. \texttt{dns.tf} declares the DNS records and \texttt{outputs.tf} exports the resulting names and the database host of each environment; these outputs are the interface with Ansible, whose inventory generator reads them in JSON form (Section~\ref{sec:impl-ansible}).

\paragraph{Policy as code.} Two Rego files in \texttt{infra/policy/}, written in the \texttt{rego.v1} dialect and sharing the package \texttt{terraform.security}, define \texttt{deny} rules evaluated against the JSON plan: \texttt{bastion\_access.rego} implements the three SSH-related checks and \texttt{droplet\_config.rego} the four droplet checks (image list, SSH key when known at plan time, bastion tag, environment tag), seven in total. Two lessons of the first version shaped the current rules. Resources of the module appear in the plan with addresses such as \texttt{module.staging.digitalocean\_firewall.frontend\_fw}, so the rules select resources by type and planned name rather than by a Terraform address or an identifier only known after the apply. The firewalls select the bastion by the tag \texttt{role:bastion-<env>}, which is why its absence is itself a violation. The earlier rule requiring \texttt{fail2ban} in the user data was removed, because fail2ban is installed by Ansible.

The rules were checked locally with Conftest against two synthetic plans written for the purpose: a compliant plan passed all seven checks and a non-compliant plan produced seven failures, one per check. No real Terraform plan was evaluated during this verification, so the behavior on the full output of the provider is not demonstrated. In the pipeline the rules are executed by Conftest in the Terraform workflow, with \texttt{--namespace terraform.security}, and by OPA in a pull-request check (Section~\ref{sec:impl-cicd}).

%========================================================================================
\section{Configuration Management with Ansible}
\label{sec:impl-ansible}

\subsection{Layout, Inventory and Playbooks}

The \texttt{ansible/} directory contains three playbooks (bootstrap, deploy and migration to the service user), twelve \texttt{roles/}, \texttt{group\_vars/} (variables common to all hosts and to the bastions, and one folder per environment with a plain \texttt{main.yml} and an encrypted \texttt{vault.yml}), \texttt{inventories/} (generated), \texttt{scripts/} (inventory generator, SSH test) and a \texttt{Makefile}. \texttt{ansible.cfg} sets the remote user to the service account \texttt{ansible}, enables host-key checking and fact caching, and enables agent forwarding and connection multiplexing for SSH; pipelining and SFTP are disabled to avoid problems with the proxied connections.

The Makefile target \texttt{inventory} exports the Terraform outputs to JSON and runs \texttt{scripts/generate\_inventory.py}, which builds the groups \texttt{staging}, \texttt{production}, \texttt{bastion}, \texttt{frontend}, \texttt{backend} and \texttt{app\_servers} from the DNS names found in the outputs and stores the database host of each environment as a variable. The bastion is contacted directly, whereas the private hosts receive an \texttt{ansible\_ssh\_common\_args} value with \texttt{-o ProxyJump} through the bastion. The target ends with a connectivity test. The connecting user is \texttt{root} until a marker file created by the bootstrap shows that the service account exists, and \texttt{ansible} afterwards.

\paragraph{Bootstrap.} Driven by \texttt{make bootstrap}, the first phase connects as \texttt{root} and applies only the \texttt{ansible\_user} role. The Makefile then creates the marker file, regenerates the inventory, and the second phase applies the remaining roles as \texttt{ansible} with privilege escalation: the hardening roles on all hosts, fail2ban on the bastion, and on the application servers \texttt{docker\_runtime} and \texttt{observability} and, on the frontend only, \texttt{nginx\_tls}. The \texttt{when} conditions rely on a fact computed in a \texttt{pre\_task}, and the play runs one host at a time by default (\texttt{serial}). \texttt{migrate-to-ansible-user.yml} performs the first phase on hosts first configured through \texttt{root}.

\paragraph{Deployment.} \texttt{deploy.yml} creates a time-stamped directory under \texttt{/opt/deployments}, logs the hosts in to GHCR with \texttt{registry\_login}, applies \texttt{app\_backend} to the backend group and \texttt{app\_frontend} to the frontend group (each with a configurable \texttt{serial}), runs \texttt{verify} on all hosts and, if a Slack webhook is defined, posts a notification. It is exposed as \texttt{make deploy ENV=staging|production} with optional \texttt{TAGS}, \texttt{LIMIT} and \texttt{DEPLOY\_VERSION}. The images \texttt{ghcr.io/oubellaismail/spring-boot-backend} and \texttt{react-frontend} are referenced with the tag \texttt{<DEPLOY\_VERSION>-<env>}, which matches the tags pushed by the CI pipeline when \texttt{DEPLOY\_VERSION} is the commit SHA. The other targets are \texttt{verify}, \texttt{ssh-test}, \texttt{lint}, \texttt{check-syntax} and \texttt{clean}.

\subsection{Roles}

Table~\ref{tab:roles-impl} lists what each role does, as read from its tasks.

\begin{table}[htbp]
\centering
\caption{Ansible roles and their actual behavior}
\label{tab:roles-impl}
\renewcommand{\arraystretch}{1.1}
\footnotesize
\begin{tabularx}{\linewidth}{@{}l X l@{}}
\toprule
\textbf{Role} & \textbf{Behavior} & \textbf{Used by} \\
\midrule
\texttt{ansible\_user} & Creates the \texttt{ansible} account in the \texttt{sudo} group, copies the authorized keys of \texttt{root}, installs a validated sudoers drop-in with passwordless sudo & bootstrap \\
\texttt{common} & Updates packages, installs common tools, sets UTC and open-files limits, creates the application system user, configures log rotation and unattended security upgrades & bootstrap \\
\texttt{ssh\_hardening} & Backs up \texttt{sshd\_config}, sets \texttt{PermitRootLogin no}, \texttt{PasswordAuthentication no}, three authentication attempts and an \texttt{AllowUsers} list; validates with \texttt{sshd -t} before restarting & bootstrap \\
\texttt{firewall} & UFW with default deny incoming; SSH from the administrator's public address (detected at run time, refused if detection fails) and the VPC range; port 80 everywhere and 443 on frontends; logging & bootstrap \\
\texttt{fail2ban} & Installs fail2ban and a \texttt{jail.local} template; adds an nginx jail on frontends; ban notification uses the e-mail ``mail with logs'' action, which assumes a local mail transfer agent & bootstrap (bastion) \\
\texttt{docker\_runtime} & Adds the Docker apt repository, installs Docker Engine and the Compose plugin, writes \texttt{daemon.json} (log rotation) & bootstrap (app servers) \\
\texttt{nginx\_tls} & Installs nginx and certbot, DH parameters, hardened TLS parameters; obtains a Let's Encrypt certificate with \texttt{certbot certonly --standalone} (nginx stopped and restarted by hooks; optionally the staging CA) and schedules renewal & bootstrap (frontend) \\
\texttt{observability} & Installs the Prometheus node exporter as a systemd service and a cron-run log-monitoring script (Slack post only if a webhook is defined) & bootstrap (app servers) \\
\texttt{registry\_login} & Logs in to GHCR with a token from the vault (\texttt{vault\_ghcr\_token}) or the environment (\texttt{GHCR\_TOKEN}); skipped when neither is defined & deploy \\
\texttt{app\_backend} & Creates the application directory, renders the environment file (mode 0600, vault values, the variable names read by Spring: \texttt{SPRING\_PROFILES\_ACTIVE=docker}, \texttt{SPRING\_DATASOURCE\_*}, \texttt{APP\_USER\_*}, \texttt{JWT\_SECRET}, \texttt{APP\_CORS\_ALLOWED\_ORIGINS}) and the Compose file, starts the container (always pulling), waits for health, runs the migration script and installs a monitoring cron job & deploy \\
\texttt{app\_frontend} & Renders the Compose file, the nginx site configuration, and cache and security-header snippets (HSTS, X-Frame-Options, a restrictive Content-Security-Policy); starts the frontend container & deploy \\
\texttt{verify} & Checks uptime, SSH/rsyslog/Docker/nginx services, application container, public health URL, TLS certificate valid for at least a week, disk usage below 90\,\% & deploy \\
\bottomrule
\end{tabularx}
\end{table}

All twelve roles are now referenced by a playbook. Four roles (\texttt{docker\_runtime}, \texttt{nginx\_tls}, \texttt{observability}, \texttt{registry\_login}) were initially left outside the playbooks and were wired into \texttt{bootstrap.yml} and \texttt{deploy.yml}, and the deployment variables were aligned with the application: the environment file uses the standard Spring datasource properties (the JDBC URL requires TLS with \texttt{sslmode=require}, as managed PostgreSQL does), activates the \texttt{docker} profile, lets Hibernate create the schema and sets the application user, signing secret and CORS origins; the health-check path is \texttt{/actuator/health}. The group variables are loaded through a symbolic link \texttt{inventories/group\_vars}, and \texttt{vault\_password\_file} was moved to the \texttt{[defaults]} section of \texttt{ansible.cfg}. These changes were validated with \texttt{ansible-playbook --syntax-check}; they have \emph{not} been executed against real hosts, so the behavior of the roles on the servers is not verified.

\subsection{Secrets and Environments}

Each environment has its own AES-256 encrypted \texttt{vault.yml} under \texttt{group\_vars/}, and the plain \texttt{main.yml} of the same folder refers to the secrets by name, for example \texttt{jwt\_secret: "\{\{ vault\_staging\_jwt\_secret \}\}"}; each vault holds the database credentials, the application user (\texttt{vault\_app\_user\_username}, \texttt{vault\_app\_user\_password}) and the signing secret. The vault password lives in the git-ignored \texttt{.vault\_pass}, which the Makefile passes to every command. The non-secret variables capture the differences between environments: domain names, GHCR image paths and tags, container resource limits, the choice between the staging and production certificate authorities, and backup and monitoring settings.

%========================================================================================
\section{CI/CD with GitHub Actions}
\label{sec:impl-cicd}

The workflows of \texttt{.github/workflows/} implement the pipeline of Figure~\ref{fig:devsecops-pipeline}: two application pipelines (backend and frontend) call reusable security workflows, one workflow handles the infrastructure, and deployment is only partly automated (Section~\ref{sec:impl-cd}). The checks and their effect on a run are summarized in Table~\ref{tab:gates}.

\paragraph{Application pipelines.} \texttt{ci-back.yaml} runs on every push to any branch that touches \texttt{back/} or the workflow files, and on pull requests to \texttt{main}, with four jobs chained by \texttt{needs}. \texttt{build} compiles the project on Java~21 with \texttt{./mvnw clean install}; \texttt{test} runs \texttt{./mvnw test} and sends the JaCoCo report to Codecov; \texttt{build-and-push-docker} logs in to Docker Hub and GHCR, builds the image with Buildx using the GitHub Actions cache and pushes it to GHCR under the commit SHA plus the suffix \texttt{-staging}; the tag is passed to the last job, which calls the reusable workflow \texttt{security-scans.yaml}. \texttt{ci-front.yaml} has the same structure: \texttt{build-and-test} sets up Node~20 with the npm cache, runs \texttt{npm ci}, builds and tests non-interactively (\texttt{npm test -{}- --watchAll=false --ci}), then the \texttt{react-frontend} image is pushed with the same tagging and \texttt{front-security-scan.yaml} is called.

\paragraph{Security workflows.} The two reusable workflows are triggered with \texttt{workflow\_call}, receive the image tag as input, and are limited to reading the repository and packages and writing security events. The backend workflow chains Gitleaks over the full history (\texttt{fetch-depth: 0}); OWASP Dependency-Check on the Maven dependencies in HTML and SARIF formats, the SARIF run using \texttt{--failOnCVSS 7}, with reports stored as artifacts and the SARIF file uploaded to code scanning; Semgrep with the \texttt{p/ci} rule set; a Trivy filesystem scan that first produces a SARIF report of HIGH and CRITICAL findings (exit code~0) and is then run again as a gate with \texttt{--exit-code 1 --severity CRITICAL --ignore-unfixed}; and a Trivy scan of the image just pushed, which uploads the SARIF report and fails on any HIGH or CRITICAL finding. The frontend workflow applies the same Gitleaks, Semgrep and Trivy steps restricted to \texttt{front/}; instead of Dependency-Check it runs \texttt{npm audit}, whose JSON report is archived and is report-only because the command is followed by \texttt{|| true}. It ends by writing a summary table to the job summary page.

\paragraph{Infrastructure workflow.} \texttt{terraform.yaml} can be started manually with the environment as a choice input, runs on pushes to \texttt{main} that modify \texttt{infra/}, and is scheduled every Monday at 06:00 UTC; a scheduled run only performs drift detection. The \texttt{plan} job installs Terraform 1.9.5 with the Terraform Cloud token, runs \textbf{tfsec} and \textbf{Checkov}, initializes the backend, selects or creates the workspace named after the environment (the name is validated against the two allowed values and passed through an environment variable instead of being interpolated into the script, which prevents shell injection), validates, and produces the plan, which is also converted to JSON for \textbf{Conftest}; the binary plan is uploaded as an artifact. The \texttt{apply} job depends on \texttt{plan}, runs only for a push or manual run on \texttt{main}, targets a GitHub \emph{environment} of the same name as the input (so that protection rules such as required reviewers can be attached), and applies exactly the saved plan, so that what was scanned is what is deployed. The \texttt{drift-detect} job selects the staging workspace and runs \texttt{terraform plan -detailed-exitcode}; exit code~2 (a difference between code and infrastructure) creates a GitHub issue from \texttt{.github/drift-template.md} with the labels \texttt{drift}, \texttt{infra}, \texttt{staging} and a priority, and the sequence is repeated for production, if its workspace exists, with a critical priority label.

A second workflow, \texttt{opa-policy-check.yaml}, is triggered by pull requests that change \texttt{infra/}. It produces a plan, renders it as JSON and evaluates \texttt{data.terraform.security.deny[\_]} with \texttt{opa eval --fail-defined}, which fails as soon as any deny rule yields a message (querying the set itself would not work, since an empty set is still a defined value). Violations are thus reported on the pull request before the merge, while the Conftest step acts as a second barrier on the main branch.

\subsection{Continuous Deployment}
\label{sec:impl-cd}

Deployment automation is the least complete part of the pipeline. \texttt{cd-back.yaml} is triggered by the completion of the workflow \texttt{Terraform Infrastructure}; it runs only if that run succeeded and was not scheduled, then opens an SSH session on a server whose address, user and key are repository secrets, updates a checkout, pulls the images and restarts the Compose services. It is a plain \texttt{docker compose up -d} on one host, with no rolling update, no automatic rollback and no call to the Ansible playbooks; because it follows the infrastructure workflow rather than the application pipeline, it is a coarse mechanism rather than a release process. \texttt{cd-front.yaml}, triggered after the frontend CI, only prints a placeholder message. The deployment path that exercises the complete configuration is \texttt{make deploy} of Section~\ref{sec:impl-ansible}, launched by an operator. Since the CI pipeline only builds images tagged with the commit SHA and the suffix \texttt{-staging}, a production deployment requires an explicit image tag.
